\documentclass[../main.tex]{subfiles}
\begin{document}
\begin{center}
    \Large{\textbf{TÓM TẮT NỘI DUNG}}\\
\end{center}
\vspace{1cm}
\indent Lãng phí thực phẩm và thiếu đói là một nghịch lý cùng tồn tại ở Việt Nam: một lượng lớn thực phẩm còn dùng được bị bỏ đi mỗi ngày trong khi nhiều cá nhân và tổ chức yếu thế vẫn thiếu bữa ăn. Hoạt động chia sẻ thực phẩm hiện nay chủ yếu mang tính tự phát và rời rạc qua mạng xã hội hoặc điện thoại, khiến người quyên góp, người nhận và lực lượng vận chuyển khó kết nối kịp thời trong khi thực phẩm lại dễ hư hỏng.

\indent Trước thực trạng đó, đồ án xây dựng Food 4 life --- một hệ thống kết nối chia sẻ thực phẩm dư thừa theo thời gian thực cho bối cảnh Việt Nam. Hệ thống được phát triển theo kiến trúc client--server, với ứng dụng di động đa nền tảng (React Native/Expo), máy chủ Node.js/Express, cơ sở dữ liệu MongoDB và kênh thời gian thực Socket.IO.

\indent Hệ thống tổ chức bốn vai trò gồm người quyên góp, người nhận, tình nguyện viên và quản trị viên; hỗ trợ luồng nghiệp vụ hai chiều khi người quyên góp đăng thực phẩm dư còn người nhận có thể chủ động đăng yêu cầu; đồng thời tích hợp gợi ý theo vị trí địa lý, điều phối tình nguyện viên bán tự động, nhắn tin, thông báo đẩy đa ngôn ngữ và theo dõi trạng thái giao nhận.

\indent Đóng góp chính của đồ án là các giải pháp kỹ thuật bảo đảm tính đúng đắn trong điều kiện vận hành không lý tưởng: ghép cặp an toàn khi tương tranh, ước lượng khoảng cách lai có dự phòng, máy trạng thái đơn tự phục hồi, mã lấy hàng chống gian lận và xác thực thu hồi tức thì. Kết quả, đồ án đã hoàn thành một sản phẩm hoàn chỉnh gồm ứng dụng di động cho ba vai trò người dùng, cổng quản trị trên web và máy chủ phụ trợ, vận hành ổn định qua các luồng nghiệp vụ chính.
\begin{flushright}
Sinh viên thực hiện\\
\begin{tabular}{@{}c@{}}
\textit{(Ký và ghi rõ họ tên)}
\end{tabular}
\end{flushright}
\end{document}